% =======================================================================
\section{Introduction}

\subsection{Background}
\label{sec:intro}

Coaches and analysts already describe team shape qualitatively: a press
that compresses space, a defensive line that holds its gap, a midfield
that loses its structure under pressure. Existing quantitative
descriptions of team shape are predominantly geometric, summarising how
players are dispersed across the pitch rather than the specific
ring-like gaps and enclosures that coaches actually describe when
discussing pressing traps or lines breaking apart.

Three such ring-like tactical structures recur in analytical and
coaching literature but lack quantitative definitions. In possession
against a compact low block, an attacking team forced to circulate the
ball around the perimeter of the defensive structure creates a
peripheral possession ring, sometimes called a horseshoe, in which
passes span from full-back to centre-back to the opposite full-back
and up the wings, leaving the central space inaccessible \cite{decroos2019}.
A pronounced peripheral ring indicates high possession volume without
penetration: the team is orbiting the opposition's defensive block
rather than entering it. In advanced positional play, a rest-defence
arc forms when deeper players spread into a shallow curved
configuration behind the front line in a 3-2-5 or 2-3-5 attacking
shape, creating a containment ring that intercepts counter-attacking
transitions and allows immediate possession recycling. In the
defensive phase, a pressing encirclement occurs when three to four
players close simultaneously around a ball-carrier following a
designated trigger pass, forming a tight ring that cuts off all exit
angles and forces a turnover in dangerous territory \cite{andrienko2017}.
These formations are described qualitatively by analysts but none has
a standard quantitative definition. Persistent homology provides one:
a ring-like gap in the spatial data produces a loop feature in the
persistence diagram, and the feature's persistence measures how stable
and geometrically significant that ring is.

A companion paper \cite{paperA} develops and validates a multi-scale
topological method that measures exactly this kind of ring-like
structure, using persistent homology to detect loops in player
positioning at two distinct spatial scales, individual player interactions and tactical-unit formations, across ten professional
matches. That validation is methodological: it establishes that the
measure is stable, reproducible, and not an artefact of parameter
choice. It does not establish what the measure means for football, or
whether it adds anything practitioners do not already have. That is the
question this paper addresses.

\subsection{Study Aims}

This study had four aims. First, to test whether topological persistence
tracks real match events in a way that is interpretable and
directionally consistent. Second, to test whether tactical-scale
persistence provides information beyond the geometric descriptors
already used in football analytics. Third, to test whether home and
away teams' tactical structures are coupled or independent within a
match. Fourth, to test whether tactical-scale persistence improves
prediction of phase of play on matches the model has not seen.

\subsection{Related Work}

Positional tracking data, from GPS or broadcast video, now underpins
much of applied performance analysis in field sports. Systematic reviews
document its role in load monitoring across football codes
\cite{tracking_review} and in tactical analysis using position data
\cite{memmert2017}. Within this literature, team shape
has typically been summarised geometrically: length, width, and
centroid distance capture how compact or stretched a team is
\cite{shape_descriptors}, and pitch-control models estimate which areas
of the field each team effectively controls at a given moment
\cite{gudmundsson2017}. Spatio-temporal pattern recognition methods
identify recurring movement sequences directly from tracking data
\cite{memmert2017}. A separate strand of work treats teams as
coordinated multi-agent systems in the spirit of collective-behaviour
research, and network-based approaches apply graph metrics to passing
structure and player interaction \cite{grund2012,buldu2018}, including
high-profile applications to identifying a team's defining tactical
characteristics.

None of these approaches directly quantifies the ring-like enclosures
, pressing traps and gaps between defensive lines, that the present
measure targets; they capture dispersion, control, or passing
structure, but not two-dimensional loop formation. The contribution of
this paper is not the topological method itself, which is developed and
validated separately \cite{paperA}, but the first systematic test of
what that method adds to this existing toolkit for football analysis
specifically.
